\chapter{Conclusões}
\label{cap:conclusoes}

\section{Considerações Finais}

Este trabalho comparou as arquiteturas monolítica e de microsserviços por meio de duas versões
funcionalmente equivalentes de uma aplicação de \textit{e-commerce}, desenvolvidas em Java com o
Spring Boot e avaliadas em um ambiente conteinerizado, com o mesmo orçamento de recursos para as
duas arquiteturas. A equivalência entre as versões foi verificada por testes automatizados, e a
comparação considerou os critérios de desempenho, escalabilidade, manutenibilidade,
complexidade de desenvolvimento e implantação, com cinco repetições de cada teste de carga e
de escalabilidade.

Os objetivos propostos foram alcançados. As diferenças conceituais entre as arquiteturas foram
apresentadas no referencial teórico, e o estado da arte reuniu trabalhos que compararam as
duas abordagens, mostrando que os seus resultados não são uniformes. As duas versões da
aplicação foram desenvolvidas com as mesmas regras de negócio, rotas e respostas, incluindo
autenticação, catálogo de produtos, personalização da loja, pedidos e pagamentos integrados ao
Mercado Pago, e foram avaliadas segundo os cinco critérios definidos.

Em resposta à questão que motivou o trabalho, os resultados indicam que, para uma aplicação de
porte pequeno a médio, desenvolvida por uma única pessoa e executada com recursos limitados, a
arquitetura monolítica se mostrou mais vantajosa na maioria dos aspectos. O monólito respondeu
mais rápido, sustentou cargas maiores com os mesmos recursos, consumiu menos processador e
memória, exigiu menos código e configuração e foi implantado com menos artefatos e em menos
tempo. Além disso, a sua replicação atrás de um balanceador de carga dobrou a vazão sem nenhuma
alteração no código, o que mostra que o monólito também pode ser escalado horizontalmente.

A arquitetura de microsserviços se mostrou vantajosa em aspectos mais específicos: o
isolamento de falhas, que manteve o catálogo funcionando enquanto o fluxo de compra estava
sobrecarregado, a possibilidade de escalar apenas os serviços mais demandados e a implantação
independente de cada serviço. Esses benefícios, porém, tiveram um custo elevado em recursos,
em código e em operação, e dependem de um conhecimento detalhado do comportamento de cada
serviço para serem aproveitados. Assim, os microsserviços tendem a compensar em contextos com
várias equipes, partes do sistema com demandas muito diferentes entre si e infraestrutura
capaz de escalar cada serviço de forma automática, condições que não estavam presentes neste
trabalho.

Os resultados também mostraram que configurações padrão e a forma de distribuir os recursos
influenciam a comparação tanto quanto a arquitetura em si. O limite do monólito foi imposto pelo
\textit{pool} de conexões com o banco de dados, e o dos microsserviços, pela fração de CPU
atribuída a cada serviço, inclusive durante o aquecimento da máquina virtual Java. Esses fatores
devem ser considerados por quem pretende escolher entre as arquiteturas ou comparar os seus
desempenhos.

\section{Contribuições}

As principais contribuições deste trabalho são:

\begin{itemize}
    \item duas implementações funcionalmente equivalentes de uma aplicação de
    \textit{e-commerce}, uma monolítica e outra em microsserviços, com integração de
    pagamentos, armazenamento de arquivos e comunicação assíncrona, cuja equivalência é
    verificada por testes automatizados;
    \item um ambiente de testes reproduzível, com orçamento de recursos equivalente para as
    duas arquiteturas, simulação do serviço de pagamento, cenários de carga, testes de
    escalabilidade vertical e horizontal e consolidação estatística automatizada dos
    resultados;
    \item uma comparação quantitativa entre as arquiteturas em cinco critérios, com resultados
    que podem auxiliar desenvolvedores na escolha da arquitetura de novos projetos;
    \item a identificação de fatores que influenciam a comparação além da própria arquitetura,
    como as configurações padrão do \textit{framework}, a divisão dos recursos entre os
    serviços e o aquecimento da máquina virtual Java com frações pequenas de processador.
\end{itemize}

\section{Trabalhos Futuros}

A partir das limitações e dos resultados deste trabalho, sugerem-se os seguintes trabalhos
futuros:

\begin{itemize}
    \item executar a comparação em um ambiente distribuído, com os serviços em máquinas
    diferentes, para medir o efeito da latência de uma rede real sobre os microsserviços;
    \item repetir os testes com as configurações ajustadas para cada arquitetura, como o
    tamanho dos \textit{pools} de conexões, e com um aquecimento mais longo da máquina virtual
    Java, para avaliar a capacidade das duas versões sem a influência das configurações
    padrão;
    \item avaliar as arquiteturas em uma plataforma de orquestração, como o Kubernetes, com
    escalabilidade automática, e comparar o custo de infraestrutura das duas versões em um
    provedor de nuvem;
    \item avaliar a resiliência das arquiteturas por meio de testes com falhas induzidas, como
    a interrupção de serviços e do banco de dados;
    \item acompanhar a evolução das duas versões ao longo do tempo e com equipes diferentes,
    para avaliar a manutenibilidade em um contexto mais próximo daquele em que os
    microsserviços são recomendados.
\end{itemize}
